\lecture{22}{20. November 2025}{Combined Loading}

\section{Combined Loadings}

\subsection{Thin-Walled Pressure Vessels}
The stresses acting in a cylindrical or spherical pressure vessel can easily be analyzed provided it has a \textit{thin wall} -- i.e. an inner-radius-to-wall--thickness ratio of 10 or more.
Specifically, when $r / t = 10$ the results of a thin-wall analysis will predict a stress that is approximately 4\% less than the actual maximum stress in the vessel.
For larger $r / t$ ratios this error will be even smaller.

\begin{figure} [ht]
  \centering
  \includegraphics[width=0.5\linewidth]{./figures/f13_1.png}
  \caption{Thin-walled cylindrical vessel.}
  \label{fig:f13_1}
\end{figure}

\subsubsection{Cylindrical vessels}
The cylindrical vessel on \Cref{fig:f13_1}a has a wall thickness $t$, inner radius $r$ and is subjected to internal gas pressure $p$.
To find the \textit{circumferential} or \textit{hoop stress} we will section the vessel by planes $a$, $b$, and $c$.
A free-body diagram of the back segment along with its contained gas is shown on \Cref{fig:f13_1}b.
For simplicity, only the loadings in the $x$ direction are shown.
These are caused bu the uniform hoop stress $\sigma_1$ acting on the vessel's wall, and the pressure acting on the vertical face of the gas.
For equilibrium in the $x$-direction we require
\begin{equation*}
  2 \left( \sigma_1 (t \, \mathrm{d}y ) \right) - p (2r \, \mathrm{d}y ) = 0 \implies \sigma_1 = \frac{pr}{t}
.\end{equation*}

The longitudinal stress can be determined by considering the left portion of section $b$ on \Cref{fig:f13_1}a.
As shown on the free-body diagram on \Cref{fig:f13_1}c, $\sigma_2$ acts uniformly throughout the wall, and $p$ acts on the section of the contained gas.
Since the mean radius is approximately equal to the vessels inner radius, equilibrium in the $y$ direction requires:
\begin{equation*}
  \sigma_2 (2\pi rt) - p \left( \pi r^2 \right) = 0 \implies \sigma_2 = \frac{pr}{2t}
.\end{equation*}
Where $\sigma_1$ and $\sigma_2$ are the normal stress in the hoop and longitudinal directions, respectively, $p$ is the internal gauge pressure, $r$ is the inner radius and $t$ the wall thickness.

By comparison, we see that the hoop stress is twice as large as the longitudinal stress.


\subsubsection{Spherical Vessels}
For a spherical vessel, the analysis is rather similar.
We section the vessel on \Cref{fig:f13_2}a in half as shown, which leaves us with the free-body diagram of \Cref{fig:f13_2}b.
Equilibrium in the $y$ direction requires
\begin{equation*}
  \omega_2 \left( 2 \pi rt \right) - p \left( \pi r^2 \right) = 0 \implies \sigma_2 = \frac{pr}{2t}
.\end{equation*}
This is the same as the longitudinal stress for a cylindrical vessel. 

\begin{figure} [ht]
  \centering
  \includegraphics[width=0.5\linewidth]{./figures/f13_2.png}
  \caption{Thin-walled spherical vessel}
  \label{fig:f13_2}
\end{figure}
